\scanpage{55}
Ens. d'orientations du cube (ou de l'espace affine qu'il engendre, ce sera provisoirement notre définition ici\ldots)
\[
\omega_C\overset{\text{déf prov}}{\simeq}\omega_E\simeq
\unsure{\omega_I|\Omega_I}\wedge\bigwedge_{i\in I}B_i
\]
L'effet de $\sigma$ sur $\omega_C$ est $(-1)^n$, donc

\uline{Prop} L'antipodisme conserve l'orientation sss $n$ pair (donc \unsure{la|le} renverse \deleted{\unsure{si|ssi}} $n$ impair)

\uline{Cor} Soit $n$ impair, $\operatorname{Aut}(C)\simeq\operatorname{Aut}^+(C)\times\{\pm1\}$

\uline{NB} (Thème de méditation géométrique) : Soit $C$ un polyèdre convexe \deleted{dans un} \inserted{engendrant l'}espace affine $E$. Alors l'ens à 2 él. $\omega_E$ peut se décrire canoniquement en termes du polyèdre combinatoire $\operatorname{Comb}(C)$.

\uline{Indication} Une orientation \inserted{$\omega$} de $E$ permet de définir \inserted{l'ens $D(\omega)$} des ``\deleted{drapeaux} \inserted{repères} directs'' $F_0\subset F_1\subset\ldots\subset F_{n-1}\subset F_n=C$ (de facettes) du polyèdre $C$, et elle est définie par la connaissance de \deleted{\unsure{ces drapeaux}} \inserted{$D(\omega)$}. Il faut donc décrire, en termes de la seule structure d'ensemble ordonné de \inserted{$\Phi=$}$\operatorname{Comb}(C)$, quelles sont les parties de \deleted{$\operatorname{Comb}(C)$} \inserted{$D(\Phi)$} (ens des \deleted{drapeaux} \inserted{repères} de $\Phi$) qui sont de la forme $D(\omega)$ [définition combinatoire de l'orientation]. \uline{Solution} Soit $D^+\subset D$, $D^-=D\setminus D^+$, on veut que \inserted{les} $\sigma_i$ $(1\leq i\leq n)$ sur $D$ transforment $D^+$ en $D^-$ (donc $D^-$ en $D^+$\ldots). En fait, il y a \unsure{exactmt|exactt|exactement} 2 telles parties : $D^+$ et $D^-$ !
% archive sheet number: 54
